\documentclass[letterpaper]{article}

\usepackage[final]{corl_2026}
\usepackage{amsmath,amssymb}
\usepackage{booktabs}
\usepackage{graphicx}
\usepackage{multirow}
\usepackage{xcolor}
\usepackage{hyperref}
\usepackage{tabularx}
\usepackage{array}

\newcolumntype{C}[1]{>{\centering\arraybackslash}p{#1}}

\title{A Systematic Diagnostic Framework for Analyzing\\Instruction Processing in Vision-Language-Action Models}

\author{
  Anonymous Authors
}

\begin{document}
\maketitle

%% ============================
%% ABSTRACT
%% ============================
\begin{abstract}
Vision-Language-Action (VLA) models fine-tune large language models (LLMs) on robotic manipulation data, yet how this training reshapes instruction processing at the token level---and whether such effects generalize across architectures---remains poorly understood.
We present a systematic diagnostic framework for token-level analysis of VLAs, combining factorial experimental design (trained/untrained $\times$ instruction presence), six complementary probing methods (RSA, CKA, R$^2$ regression, counterfactual sensitivity, attention analysis, classification probes), and cross-architecture validation.
We apply this framework to three VLA families---OpenVLA (Llama-2-7B), TracVLA (Phi-3-mini), and NORA (Qwen2.5-VL-3B)---across four LIBERO task suites, yielding a 10-cell evaluation matrix.
The framework reveals two classes of effects: \emph{training-induced} phenomena (R$^2$ exit collapse where action-relevant information drops at exit layers, trained $\Delta$=0.04--0.69; and heightened counterfactual sensitivity, trained 2--10$\times$ untrained) that appear directionally consistent across architectures, and \emph{base-model-dependent} properties (L1 attention head counts) that reflect pre-existing LLM structure rather than VLA training.
We further document metric dissociation: RSA and CKA yield reversed ordering of representational similarity, underscoring why multi-metric frameworks are necessary.
Our framework is released as a reusable toolkit for VLA representational analysis.
\end{abstract}

%% ============================
%% 1. INTRODUCTION
%% ============================
\section{Introduction}

Vision-Language-Action (VLA) models \cite{brohan2023rt2,kim2024openvla} integrate vision encoders with large language models (LLMs) to produce robotic actions conditioned on natural language instructions and visual observations.
Despite growing deployment, the internal mechanisms by which these models process instructions and route information from language tokens to action outputs remain largely opaque.
Understanding these mechanisms is critical: if instruction processing follows predictable patterns, architectural design and debugging can be guided by diagnostic analysis rather than trial-and-error.

Recent behavioral studies have documented instruction underutilization and visual dominance in VLAs, but these operate at the input-output level and do not decompose the representational roles of different token positions across layers.
The key open questions are: (1) What does VLA fine-tuning actually change in the LLM's internal representations? (2) Are these changes architecture-specific artifacts or general properties of VLA training? (3) Can we build reusable diagnostic tools that answer such questions for any decoder-only VLA?

We address these questions by developing a \textbf{systematic diagnostic framework} for token-level VLA analysis.
Our framework combines: (a) factorial experimental design that disentangles training effects from instruction effects; (b) six complementary probing methods that triangulate findings across geometric, functional, and attentional perspectives; and (c) a cross-architecture validation protocol applied to three VLA families with distinct LLM backbones.

Applying this framework across three architectures (OpenVLA/Llama-2-7B, TracVLA/Phi-3-mini, NORA/Qwen2.5-VL-3B) and four LIBERO task suites, we find:

\begin{itemize}
    \item \textbf{Framework contribution}: A reusable diagnostic toolkit validated on 10 architecture$\times$task cells, capable of decomposing encoder vs.\ LLM contributions, distinguishing training-induced from base-model effects, and revealing metric-dependent conclusions.
    \item \textbf{Training-induced effects}: R$^2$ exit collapse (9/10 cells show trained $>$ untrained delta) and counterfactual sensitivity (11/12 settings show trained $\gg$ untrained) are directionally consistent across architectures, though magnitudes vary substantially (exit $\Delta$: 0.004--0.691).
    \item \textbf{Base-model attribution}: L1 instruction-dominant head counts are determined by the base model's pre-training, not VLA fine-tuning---a finding that our framework's controlled comparison uniquely reveals.
    \item \textbf{Metric dissociation}: RSA and CKA yield reversed representational similarity orderings, demonstrating that single-metric analyses can produce misleading conclusions about VLA representations.
\end{itemize}

We release the framework as a diagnostic toolkit and discuss its limitations, including the absence of behavioral validation and the need for NLP fine-tuning controls to establish VLA-specificity of exit collapse.

%% ============================
%% 2. RELATED WORK
%% ============================
\section{Related Work}

\paragraph{Vision-Language-Action Models.}
VLAs extend vision-language models (VLMs) with action prediction capabilities.
RT-2 \cite{brohan2023rt2} demonstrated that LLM-scale VLMs could generate robotic actions as text tokens.
OpenVLA \cite{kim2024openvla} open-sourced a 7B-parameter VLA trained on the Open X-Embodiment dataset.
More recent architectures include TracVLA (Phi-3-mini backbone) and NORA (Qwen2.5-VL backbone), which employ different vision encoders and LLM architectures.
Our work does not propose a new VLA but provides diagnostic tools to analyze how existing VLAs process instructions internally.

\paragraph{Mechanistic Analysis of Language Models.}
Probing classifiers \cite{belinkov2017neural,hewitt2019structural}, representational similarity analysis (RSA) \cite{kriegeskorte2008representational}, and centered kernel alignment (CKA) \cite{kornblith2019similarity} are standard tools for analyzing neural network representations.
Attention-based analyses have revealed specialized heads in LLMs for specific linguistic functions.
We adapt these methods to the multi-modal VLA setting, where image, instruction, and action tokens share a single transformer backbone.

\paragraph{Behavioral Analysis of VLAs.}
Grant et al.\ (2026) and Haon et al.\ (2025) studied VLA behavior through input perturbation, finding visual dominance and instruction underutilization.
These studies operate at the input-output level.
Our framework complements them by providing layer-by-layer internal analysis, disentangling what training changes from what the base model already provides.

%% ============================
%% 3. METHOD: DIAGNOSTIC FRAMEWORK
%% ============================
\section{Diagnostic Framework}

We describe a systematic framework for analyzing instruction processing in decoder-only VLAs.
The framework consists of three components: factorial experimental design (\S\ref{sec:factorial}), a multi-method analysis suite (\S\ref{sec:methods}), and a cross-architecture validation protocol (\S\ref{sec:protocol}).

\subsection{Factorial Experimental Design}
\label{sec:factorial}

To disentangle training effects from architectural properties, we compare each VLA under multiple conditions:

\begin{itemize}
    \item \textbf{Trained}: VLA fine-tuned on task-specific data (e.g., LIBERO suite).
    \item \textbf{Untrained}: Same architecture with either (a) random initialization (OpenVLA), (b) base VLM weights without action training (TracVLA/Phi-3-Vision, NORA/Qwen2.5-VL), or (c) general-purpose OXE-pretrained weights (openvla-7b).
    \item \textbf{Instruction presence}: With original task instruction vs.\ without instruction (empty/shuffled/random).
\end{itemize}

This $2 \times 2$ factorial design (training $\times$ instruction) enables controlled attribution of observed effects.
We additionally validate that the vision encoder (SigLIP) and projector remain frozen during training (cosine similarity = 1.0 between trained and untrained projector weights), confirming that observed representation changes originate in the LLM backbone.

\subsection{Multi-Method Analysis Suite}
\label{sec:methods}

We employ six complementary probing methods, each capturing different aspects of representation:

\begin{enumerate}
    \item \textbf{Representational Similarity Analysis (RSA)}: Measures geometric similarity between model representations and a reference (action space). Captures global geometric structure.
    \item \textbf{Centered Kernel Alignment (CKA)}: Measures global structural alignment. Provides a complementary view to RSA---as we show, the two can yield divergent conclusions.
    \item \textbf{R$^2$ regression probes}: Linear probes predicting action variables from per-layer representations. Measures functional (decodable) information content.
    \item \textbf{Counterfactual sensitivity}: Measures representation shift when instructions are replaced with counterfactual alternatives (empty, shuffled, random, semantically different). Quantifies instruction dependence.
    \item \textbf{Attention analysis}: Per-head and per-layer attention fraction allocated to instruction vs.\ image tokens. Identifies instruction-dominant attention heads.
    \item \textbf{Classification probes}: Cross-instruction $k$-way classification (same image, $k$ different instructions) and subtask probes. Measures instruction discriminability.
\end{enumerate}

\subsection{Cross-Architecture Validation Protocol}
\label{sec:protocol}

We apply the framework to three VLA families spanning diverse LLM backbones:

\begin{itemize}
    \item \textbf{OpenVLA}: Llama-2-7B (32 layers, multi-head attention), DinoSigLIP vision encoder.
    \item \textbf{TracVLA}: Phi-3-mini (32 layers, multi-head attention), Phi-3-Vision encoder.
    \item \textbf{NORA}: Qwen2.5-VL-3B (36 layers, grouped-query attention with 16 heads / 2 KV heads), InternViT encoder.
\end{itemize}

Each is evaluated on up to four LIBERO task suites (Goal, Spatial, Object, Long), yielding a maximum of 12 cells.
NORA Object and Long checkpoints were unavailable, resulting in 10/12 filled cells.

%% ============================
%% 4. EXPERIMENTS
%% ============================
\section{Experiments}

We organize results around the diagnostic capabilities our framework demonstrates: encoder--LLM decomposition (\S\ref{sec:encoder}), cross-architecture exit collapse analysis (\S\ref{sec:exit}), counterfactual sensitivity (\S\ref{sec:cf}), training vs.\ base-model attribution (\S\ref{sec:d090}), preliminary observations on a routing gradient (\S\ref{sec:gradient}), metric dissociation (\S\ref{sec:metric}), and encoder comparison (\S\ref{sec:encodercomp}).

% ------- 4.1 -------
\subsection{Encoder--LLM Decomposition}
\label{sec:encoder}

Our framework decomposes VLA representations into encoder and LLM contributions.
On OpenVLA, the SigLIP+projector yields RSA = 0.615 with action-space geometry, while the full VLA (trained) peaks at RSA = 0.488 and the untrained VLA at RSA = 0.485.
The $\Delta = 0.127$ degradation from projector to LLM output quantifies the geometric cost of LLM processing.
DINO (standalone vision model, no language) serves as a reference at RSA = 0.388.

\paragraph{Per-layer RSA trajectory.}
RSA follows an inverted-U pattern in the LLM: from projector output (0.615), it drops to the first LLM layer (drop of 0.13), then peaks at layer 9 (trained: 0.488) or layer 12 (untrained: 0.485), before declining through deeper layers.
This pattern is confirmed only on OpenVLA (Llama-2-7B); TracVLA and NORA per-layer RSA curves were not computed, so generality of the inverted-U cannot be claimed.

\paragraph{Factorial instruction invariance.}
The $2 \times 2$ factorial design (training $\times$ instruction) yields RSA values in a narrow range [0.480, 0.489], with training effect $\Delta = 0.003$ and instruction effect $\Delta = 0.004$ (Table~\ref{tab:2x2}).
Bootstrap confidence intervals include zero for both effects.
This is structurally expected: image tokens precede instruction tokens in the causal mask and thus cannot attend to them, so image-token RSA is architecturally invariant to instruction content.

\begin{table}[t]
\centering
\caption{$2 \times 2$ factorial RSA analysis on OpenVLA. All conditions yield near-identical RSA, consistent with causal mask structure guaranteeing instruction invariance of image token representations.}
\label{tab:2x2}
\small
\begin{tabular}{lcc}
\toprule
\textbf{Condition} & \textbf{RSA} & \textbf{Reviewer interpretation} \\
\midrule
Trained + Instruction & 0.488 & Frozen encoder dominates \\
Trained -- Instruction & 0.489 & Image RSA independent of text \\
Untrained + Instruction & 0.485 & Same encoder, same RSA \\
Untrained -- Instruction & 0.480 & Null result confirmed \\
\bottomrule
\end{tabular}
\end{table}

We further verify instruction invariance with four instruction conditions (normal, random, shuffled, empty), finding maximum absolute $\Delta = 0.017$ (Table~\ref{tab:instcontrol}).
This validates the framework's sanity-checking capability but does not constitute a novel finding.

\begin{table}[t]
\centering
\caption{Instruction invariance control on OpenVLA. RSA is stable across all instruction conditions ($\Delta_{\max} = 0.017$), confirming causal mask guarantees.}
\label{tab:instcontrol}
\small
\begin{tabular}{lc}
\toprule
\textbf{Instruction condition} & \textbf{RSA} \\
\midrule
Normal & 0.489 \\
Random & 0.484 \\
Shuffled & 0.483 \\
Empty & 0.472 \\
\bottomrule
\end{tabular}
\end{table}

\paragraph{Six-method consistency.}
Across six probing methods and two benchmarks, the trained--untrained delta remains small: RSA $\Delta = 0.003$, R$^2$ $\Delta = 0.015$, CKA $\Delta = -0.006$, action probe $\Delta = 0.045$, delta probe $\Delta = 0.011$, classification $\Delta = 0.001$ (Table~\ref{tab:sixmethod}).
All methods consistently show trained $\approx$ untrained for distributed representations, as expected given the frozen encoder.
The value of this analysis lies in demonstrating the framework's multi-method triangulation capability, not in the finding itself.

\begin{table}[t]
\centering
\caption{Six-method consistency for trained vs.\ untrained representations on OpenVLA. All deltas are small, confirming frozen-encoder dominance across diverse metrics.}
\label{tab:sixmethod}
\small
\begin{tabular}{lcl}
\toprule
\textbf{Method} & \textbf{Trained--Untrained $\Delta$} & \textbf{Interpretation} \\
\midrule
RSA & 0.003 & Frozen encoder dominates \\
R$^2$ probe & 0.015 & Similar L0 inputs \\
CKA & $-0.006$ & Reversal (see \S\ref{sec:metric}) \\
Action probe & 0.045 & Marginal routing signal \\
Delta probe & 0.011 & Similar input structure \\
Classification & 0.001 & Trivially easy for both \\
\bottomrule
\end{tabular}
\end{table}

% ------- 4.2 -------
\subsection{R$^2$ Exit Collapse Across Architectures}
\label{sec:exit}

The most prominent training-induced effect our framework detects is \emph{R$^2$ exit collapse}: action-relevant information (measured by linear R$^2$ probes on the last pre-action token) peaks at intermediate layers but drops sharply at exit layers in trained models, while untrained models show a smaller or absent decline.

\paragraph{Full-layer R$^2$ scan on OpenVLA.}
A 33-layer $\times$ 2-condition $\times$ 2-aggregation scan (132 probes) reveals: trained model peaks at R$^2$ = 0.678 (layer 13) and drops to R$^2$ = 0.422 at exit (layer 32), yielding $\Delta = 0.256$.
The untrained model (random init) peaks at R$^2$ = 0.609 (layer 32) with no exit decline.
The Llama-2 base model (pretrained LLM, no VLA training) shows a milder decline: peak R$^2$ = 0.645 (layer 9), exit R$^2$ = 0.555, $\Delta = 0.090$ (Table~\ref{tab:r2exit}).

A critical missing control is an NLP-fine-tuned LLM: if Llama-2 fine-tuned on any supervised task shows comparable exit collapse, the effect would be a generic fine-tuning property rather than VLA-specific.
We flag this as a key limitation.

\begin{table}[t]
\centering
\caption{R$^2$ exit collapse on OpenVLA (full-layer scan). Trained model shows the largest peak-to-exit delta. The critical NLP fine-tuning control is missing.}
\label{tab:r2exit}
\small
\begin{tabular}{lccl}
\toprule
\textbf{Condition} & \textbf{Peak R$^2$ (layer)} & \textbf{Exit R$^2$} & \textbf{$\Delta$} \\
\midrule
Trained & 0.678 (L13) & 0.422 & 0.256 \\
Untrained (random) & 0.609 (L32) & n/a & n/a \\
Llama-2 base & 0.645 (L9) & 0.555 & 0.090 \\
NLP fine-tuned & ? & ? & ? (missing) \\
\bottomrule
\end{tabular}
\end{table}

\paragraph{Cross-architecture matrix.}
Table~\ref{tab:matrix} presents the 3-architecture $\times$ 4-suite R$^2$ exit collapse matrix.
Of 10 filled cells, 9 show trained $\Delta >$ untrained $\Delta$ (directional consistency), but magnitudes range from 0.004 to 0.691---a 172$\times$ range.
Six cells show clear collapse ($\Delta > 0.10$), while four are weak or absent.

\begin{table}[t]
\centering
\caption{Cross-architecture R$^2$ exit collapse matrix. Trained delta = peak R$^2$ minus exit R$^2$. Ratio = trained $\Delta$ / untrained $\Delta$. 9/10 cells are directionally consistent (trained $>$ untrained), but magnitude varies 172$\times$.}
\label{tab:matrix}
\small
\setlength{\tabcolsep}{3pt}
\begin{tabular}{llcccl}
\toprule
\textbf{Architecture} & \textbf{Suite} & \textbf{Tr.\ $\Delta$} & \textbf{Untr.\ $\Delta$} & \textbf{Ratio} & \textbf{Verdict} \\
\midrule
OpenVLA & Goal & 0.069 & 0.015 & 4.6$\times$ & Moderate \\
OpenVLA & Spatial & 0.040 & 0.023 & 1.7$\times$ & Weak \\
OpenVLA & Object & 0.121 & 0.076 & 1.6$\times$ & Moderate \\
OpenVLA & Long (OFT) & 0.004 & 0.016 & 0.25$\times$ & Absent \\
\midrule
TracVLA & Goal & 0.090 & 0.048 & 1.9$\times$ & Weak \\
TracVLA & Spatial & 0.138 & 0.016 & 8.6$\times$ & Strong \\
TracVLA & Object & 0.294 & 0.009 & 32.7$\times$ & Strong \\
TracVLA & Long & 0.198 & 0.028 & 7.1$\times$ & Strong \\
\midrule
NORA & Goal & 0.691 & 0.044 & 15.7$\times$ & Extreme \\
NORA & Spatial & 0.302 & 0.037 & 8.2$\times$ & Strong \\
\bottomrule
\end{tabular}
\end{table}

Notable observations:
(1) OpenVLA (the only community-released checkpoint) shows the \emph{weakest} collapse, with 2/4 suites weak or absent.
(2) NORA Goal ($\Delta = 0.691$) is an order of magnitude above OpenVLA, raising questions about whether strong collapse correlates with self-trained checkpoints.
(3) The peak layer standard deviation across architectures is 6.2 layers, indicating that \emph{where} collapse occurs is also architecture-dependent.

\paragraph{Cross-benchmark consistency within OpenVLA.}
The standardized pipeline re-run on OpenVLA Goal yields exit $\Delta = 0.070$ (vs.\ original 0.069), confirming reproducibility.
OpenVLA Spatial yields $\Delta = 0.040$---substantially weaker than Object (0.121) and Goal (0.070), suggesting exit collapse intensity is task-dependent.
Table~\ref{tab:openvla_suites} summarizes OpenVLA cross-suite results.

\begin{table}[t]
\centering
\caption{OpenVLA cross-suite exit collapse. Mean $\Delta = 0.059$, the weakest across all three architectures.}
\label{tab:openvla_suites}
\small
\begin{tabular}{lccl}
\toprule
\textbf{Suite} & \textbf{Trained $\Delta$} & \textbf{Standardized} & \textbf{Verdict} \\
\midrule
Goal & 0.069 & 0.070 & Moderate \\
Spatial & 0.040 & 0.040 & Weak \\
Object & 0.121 & --- & Moderate-strong \\
Long (OFT) & 0.004 & --- & Absent \\
\bottomrule
\end{tabular}
\end{table}

\paragraph{LIBERO-Spatial cross-benchmark detail.}
On LIBERO-Spatial (the most spatially demanding suite), OpenVLA trained image-mean peak R$^2$ = 0.696 and last-preaction peak R$^2$ = 0.675.
The trained exit $\Delta$ = 0.040 ($p = 0.001$ via bootstrap) is statistically significant but practically small.
Cross-benchmark consistency is partial: the effect direction holds but magnitude is task-dependent.

\paragraph{Matrix coverage and gaps.}
Two cells are missing (NORA Object/Long), yielding 10/12 coverage.
Of the 10 filled cells, 6 show clear collapse ($\Delta > 0.10$), and 4 are weak or absent.
The high variability in both magnitude and peak layer location means exit collapse is better characterized as a ``directionally consistent but magnitude-variable observation'' than an ``architecture-independent signature.''

% ------- 4.3 -------
\subsection{Counterfactual Sensitivity Analysis}
\label{sec:cf}

Counterfactual sensitivity measures how much exit-layer representations shift when instructions are replaced.
Across all tested conditions, trained models show substantially higher sensitivity than untrained models (Table~\ref{tab:cf}).

\begin{table}[t]
\centering
\caption{Counterfactual sensitivity across architectures. Trained models consistently show higher sensitivity (2--10$\times$), but the untrained baselines are not functional models, limiting the fairness of comparison.}
\label{tab:cf}
\small
\setlength{\tabcolsep}{3pt}
\begin{tabular}{llccl}
\toprule
\textbf{Arch} & \textbf{Suite} & \textbf{Trained $d$} & \textbf{Untrained} & \textbf{Ratio} \\
\midrule
OpenVLA & Object & 1.21 & 0.569 & 2.1$\times$ \\
TracVLA & Goal & 1.52 & 0.153 & 9.9$\times$ \\
TracVLA & Spatial & 1.33 & 0.180 & 7.4$\times$ \\
TracVLA & Object & 1.21 & 0.185 & 6.5$\times$ \\
TracVLA & Long & 1.47 & 0.173 & 8.5$\times$ \\
NORA & Goal & 1.0 (sat.) & 0.081 & 12.3$\times$ \\
\bottomrule
\end{tabular}
\end{table}

In 11/12 tested settings, ``empty'' instruction produces the largest representation shift in trained models, suggesting that instruction removal is the most disruptive counterfactual.
One exception is OpenVLA Long, where ``shuffled'' exceeds ``empty.''

\paragraph{Limitations of counterfactual analysis.}
The trained vs.\ untrained comparison has a fairness concern: untrained models are either randomly initialized (OpenVLA) or base VLMs without action heads (TracVLA, NORA)---they are not functional action-producing models.
A fairer comparison would use an NLP-fine-tuned LLM, which processes instructions but was not trained on robotic data.
Additionally, NORA Goal yields a saturated cosine distance of 1.0, which may reflect an architectural artifact (exit-layer representations becoming orthogonal across conditions) rather than genuine routing.

% ------- 4.4 -------
\subsection{Distinguishing Training-Induced from Base-Model Effects}
\label{sec:d090}

A central capability of our framework is separating what VLA training creates from what the base model already provides.
This is demonstrated most clearly through L1 attention head analysis across three architectures (Table~\ref{tab:d090}).

\begin{table}[t]
\centering
\caption{L1 instruction-dominant heads: base-model-dependent, not training-induced. TracVLA untrained (Phi-3-Vision) already has 32/32 heads instruction-dominant; training \emph{reduces} this count. NORA is unchanged by training.}
\label{tab:d090}
\small
\begin{tabular}{llccl}
\toprule
\textbf{Arch} & \textbf{Backbone} & \textbf{Trained} & \textbf{Untrained} & \textbf{Training effect} \\
\midrule
OpenVLA & Llama-2-7B & 24/32 & 0/32 (rand.) & Artifact: rand.\ has no LM \\
OpenVLA & OXE-pretr. & 24/32 & 25/32 & Zero \\
TracVLA & Phi-3-mini & 27--30/32 & 32/32 (VLM) & Negative \\
NORA & Qwen2.5-VL & 12--13/16 & 13--14/16 & Zero \\
\bottomrule
\end{tabular}
\end{table}

The OpenVLA trained-vs-random comparison (0 $\rightarrow$ 24 heads) initially appeared to demonstrate that training creates instruction routing.
However, comparing with the OXE-pretrained checkpoint (25/32 heads, \emph{no task-specific training}) reveals that general LLM/VLM pretraining already establishes these attention patterns.
TracVLA's untrained Phi-3-Vision backbone has \emph{all} 32 heads instruction-dominant; training actually reduces this to 27--30.
NORA shows no change (13/16 $\rightarrow$ 12--13/16).

This result---which removes one of the most visually striking findings from the contribution list---demonstrates the framework's value: without controlled cross-condition comparison, the L1 gradient would have been misattributed to VLA training.

\paragraph{Contribution audit.}
After this analysis, the remaining training-induced effects are: (1) R$^2$ exit collapse (6/10 cells clear, VLA-specificity questionable without NLP control) and (2) counterfactual sensitivity (11/12 settings, but unfair baseline).
Table~\ref{tab:audit} summarizes the post-analysis contribution status.

\begin{table}[t]
\centering
\caption{Contribution audit after distinguishing training-induced from base-model effects.}
\label{tab:audit}
\small
\begin{tabular}{lll}
\toprule
\textbf{Claimed contribution} & \textbf{Status} & \textbf{VLA-specific?} \\
\midrule
L1 routing gradient & Removed (base-model) & No \\
R$^2$ exit collapse & Retained (6/10 clear) & Needs NLP control \\
CF sensitivity & Retained (11/12) & Unfair baseline \\
Three-tier gradient & Thin (N=1 at Tier-2) & Preliminary \\
Diagnostic framework & Intact & Yes \\
Metric dissociation & Intact & Methodological \\
\bottomrule
\end{tabular}
\end{table}

% ------- 4.5 -------
\subsection{Preliminary Observation: Training-Stage Gradient}
\label{sec:gradient}

Our framework enables comparison across training stages.
We observe a preliminary three-tier pattern (Table~\ref{tab:gradient}):

\begin{table}[t]
\centering
\caption{Three-tier training-stage comparison. Tier-2 has only N=1 data point (OpenVLA OXE-pretrained), making ``gradient'' a preliminary characterization rather than a confirmed finding.}
\label{tab:gradient}
\small
\setlength{\tabcolsep}{3pt}
\begin{tabular}{lcccc}
\toprule
\textbf{Tier} & \textbf{Description} & \textbf{N} & \textbf{Exit collapse?} & \textbf{CF sensitivity} \\
\midrule
1: VLM base & No action training & 3 & No & $<0.2$ \\
2: OXE-pretr. & General action & 1 & No & $d = 1.048$ \\
3: Task-finetuned & Specific action & $\sim$10 & Yes (6/10) & $d = 1.2$--$1.5$ \\
\bottomrule
\end{tabular}
\end{table}

Tier-2 has only one data point (openvla-7b pretrained on 970K OXE episodes), so the ``gradient'' is better described as a step function: VLM (no routing) $\rightarrow$ VLA (routing).
Multiple Tier-2 checkpoints (different pretraining datasets, different architectures) would be needed to confirm a true gradient.
We present this as a preliminary observation that our framework is designed to investigate, not as a confirmed finding.

% ------- 4.6 -------
\subsection{Metric Dissociation: RSA vs.\ CKA}
\label{sec:metric}

A methodologically important finding is that RSA and CKA produce \emph{reversed} orderings of representational similarity (Table~\ref{tab:metric}).

\begin{table}[t]
\centering
\caption{RSA vs.\ CKA ordering reversal. The two metrics yield contradictory rankings, demonstrating the necessity of multi-metric frameworks.}
\label{tab:metric}
\small
\begin{tabular}{lll}
\toprule
\textbf{Metric} & \textbf{Ordering} & \textbf{Implication} \\
\midrule
RSA (geometric) & SigLIP $>$ T $>$ U $>$ DINO & Training degrades geometry \\
CKA (global) & U $>$ T $>$ SigLIP $>$ DINO & Training preserves structure \\
R$^2$ (functional) & T$_{\text{readout}}$ $>$ DINO $\approx$ Proj & Training enables decoding \\
\bottomrule
\end{tabular}
\end{table}

Under RSA, training \emph{degrades} representational geometry relative to the encoder output.
Under CKA, the untrained model \emph{better preserves} global structure than the trained one.
Under R$^2$ probes, training \emph{enables} action decoding at readout positions.
The linear CKA at the last pre-action layer is 0.127, far below RSA's peak of 0.488 at the same position.

This dissociation has two implications.
First, claims about VLA representations are metric-dependent: what looks like ``degradation'' under one metric appears as ``preservation'' under another.
Second, multi-metric frameworks are necessary---single-metric analyses risk producing metric-artifact conclusions.
However, this same dissociation undermines the authority of \emph{any} specific quantitative finding in this paper, including R$^2$ exit collapse.

% ------- 4.7 -------
\subsection{Vision Encoder Comparison}
\label{sec:encodercomp}

The framework's encoder decomposition capability enables cross-encoder comparison.
Table~\ref{tab:encoder} presents RSA and R$^2$ for three encoder types.

\begin{table}[t]
\centering
\caption{Encoder comparison (under RSA and R$^2$). DINO (no language supervision) outperforms both vision-language encoders, raising questions about the practical relevance of instruction routing.}
\label{tab:encoder}
\small
\begin{tabular}{lccl}
\toprule
\textbf{Encoder} & \textbf{RSA} & \textbf{R$^2$} & \textbf{Language?} \\
\midrule
DINO cls & 0.692 & 0.634 & No \\
SigLIP+Proj & 0.615 & 0.630 & Yes \\
CLIP ViT-L & 0.557 & 0.587 & Yes \\
\bottomrule
\end{tabular}
\end{table}

DINO (self-supervised, no language pretraining) achieves the highest RSA (0.692) and R$^2$ (0.634), outperforming both vision-language encoders.
CLIP ViT-L ranks last despite vision-language pretraining.
This suggests that vision-language pretraining does not inherently produce better action-relevant representations than self-supervised learning, though this comparison is limited to the RSA metric and single-architecture setting.

On OpenVLA's DinoSigLIP fused backbone, branch-level analysis shows near-identical RSA for the DINOv2 branch (0.619) and SigLIP branch (0.621), with fusion yielding a modest improvement (0.634).
Post-projector RSA drops to 0.615, indicating minor information loss at the projection layer.

%% ============================
%% 5. DISCUSSION
%% ============================
\section{Discussion}

\paragraph{Framework as primary contribution.}
The main contribution of this work is the diagnostic framework itself: its factorial design, multi-method suite, and cross-architecture protocol enable systematic analysis of any decoder-only VLA.
The findings we report---exit collapse, counterfactual sensitivity, base-model attribution---are demonstrations of what the framework reveals, not standalone claims of architectural universality.

\paragraph{Honest assessment of exit collapse.}
R$^2$ exit collapse is directionally consistent across architectures (9/10 cells) but not magnitude-consistent (0.004--0.691).
The community checkpoint (OpenVLA) shows the weakest effects, raising the possibility that strong collapse is partially a training-procedure artifact.
Without an NLP-fine-tuning control, we cannot establish that exit collapse is VLA-specific rather than a generic property of supervised fine-tuning on decoder-only models.

\paragraph{Metric dependence.}
The RSA--CKA dissociation is both a limitation and a contribution.
It limits the strength of any single finding in this paper.
It motivates the multi-metric approach and suggests that the VLA analysis community should standardize metric reporting, analogous to reporting both accuracy and F1 in NLP.

\paragraph{Absence of behavioral validation.}
This work is entirely representational---we do not demonstrate that the patterns we observe (exit collapse, attention routing) causally influence robot task performance.
For a robotics venue, this is a significant limitation.
Future work should connect representational diagnostics to behavioral outcomes through intervention experiments (e.g., early exit, attention masking).

\paragraph{Limitations summary.}
(1) No NLP fine-tuning control for exit collapse VLA-specificity.
(2) Metric dissociation undermines quantitative authority.
(3) 4/10 matrix cells show weak or absent collapse.
(4) NORA Object/Long missing.
(5) Three-tier gradient has N=1 at Tier-2.
(6) No behavioral validation.
(7) Cross-architecture RSA curves not computed for TracVLA/NORA.
(8) Encoder comparison limited to OpenVLA's architecture.

%% ============================
%% 6. CONCLUSION
%% ============================
\section{Conclusion}

We introduced a systematic diagnostic framework for analyzing instruction processing in Vision-Language-Action models.
Applied to three VLA architectures across four task suites, the framework reveals training-induced effects (R$^2$ exit collapse, counterfactual sensitivity) that are directionally consistent but magnitude-variable across architectures, and base-model properties (L1 attention heads) that were previously misattributed to VLA training.
The framework also uncovers metric dissociation between RSA and CKA, demonstrating why multi-metric analysis is necessary.

Our primary contribution is the framework itself---a reusable toolkit for VLA representational analysis---rather than any single finding about routing.
The findings we present are first demonstrations of the framework's capabilities, subject to the limitations we have discussed (missing NLP controls, no behavioral validation, variable effect magnitudes).
We believe this framework provides a foundation for more targeted, better-controlled investigations of how VLAs process multimodal information.

\bibliography{references}
\bibliographystyle{plainnat}

\end{document}
